% ch04-systems.tex -- RV32I assembly II: functions, the calling
% convention, stack frames, recursion, and reading real compiler output.
\chapter{Functions and the Stack}
\label{ch:stack}

A function call looks like one line of source and feels free. On the
machine it is two instructions of mechanism --- \texttt{call} writes
the return address into \texttt{ra} and jumps; \texttt{ret} jumps back
--- wrapped in a \emph{convention}: an agreement about which registers
survive the trip and where memory holds the values that do not fit in
registers. In the last chapter every program was a leaf: it never
called anything, so \texttt{ra} stayed safe in a register and the
stack was never touched. This chapter breaks that. We give the machine
a real \texttt{sp}, write a recursive \texttt{fib} that saves and
restores \texttt{ra} correctly, measure what a call costs in retired
instructions, and then look at real compiler output --- where the
calling convention lives inside every function you will ever write.
By the end, \texttt{addi sp,sp,-16 / sw ra,12(sp) / sw s0,8(sp)} reads
like a signature, not a spell.

\begin{wins}[What you will be able to do after this chapter]
You will be able to read a function's prologue and epilogue in real
compiler output and name every instruction's job --- why \texttt{sp}
moves, why \texttt{ra} and the saved registers go where they go. You
will be able to write a recursive function in RV32I that saves and
restores state correctly, and to reason about what a call costs, in
instructions, compared with the work inside it.
\end{wins}

\begin{rulebox}[Before you start]
Two things carry over; each is refreshed in a few lines.
\emph{RV32I} (ch03 of this book): the register file, the ABI names,
loads, stores and branches --- this chapter is written in that
language, so keep the names \texttt{ra}, \texttt{sp}, \texttt{a0}
and \texttt{s0}--\texttt{s11} within reach.
\emph{A stack} (Book 1, ch04): push on top, pop from top, last in
first out. The call stack is the same idea with a twist: each call
pushes a \emph{frame} --- a block holding the return address, saved
registers and locals --- and the pointer \texttt{sp} is just the
top-of-stack marker that moves as frames come and go.
\end{rulebox}

\section{The contract in one paragraph}
\label{sec:contract04}
\index{calling convention}\index{psABI}\index{register!caller-saved}
\index{register!callee-saved}

The RISC-V calling convention (the \emph{psABI}) divides the 32
registers into two families \emph{by agreement}, not by hardware.
\textbf{Caller-saved} registers (\texttt{t0}--\texttt{t6},
\texttt{a0}--\texttt{a7}, \texttt{ra}) may be destroyed by any call:
if the caller needs the value afterwards, the \emph{caller} saves it
before the call. \textbf{Callee-saved} registers (\texttt{s0}--\texttt{s11},
\texttt{sp}) must be unchanged when the callee returns: if the callee
wants to use them, the \emph{callee} saves and restores them. The
arguments go in \texttt{a0}--\texttt{a7}, the return value comes back
in \texttt{a0}, and \texttt{sp} must stay 16-byte aligned at every
call site. That is the whole contract. Everything else --- the frame,
the spill slots, the order of saves --- is what the contract costs
when calls nest.

\section{The stack frame}
\label{sec:frame04}
\index{stack frame}\index{sp@\texttt{sp}}\index{ra@\texttt{ra}}

When a function calls another function, its own \texttt{ra} is about
to be overwritten --- \texttt{call} always writes the same register.
So before a non-leaf function calls anything, it builds a
\emph{stack frame}: a fixed-size block of memory below \texttt{sp}
that holds what the function cannot afford to lose. Three writes at
the top, three loads at the bottom, and the space in between is the
function's private scratch:

\begin{lstlisting}
fib:
        addi  sp, sp, -16    # carve 16 B below sp (aligned)
        sw    ra, 12(sp)     # save MY return address
        sw    s0, 8(sp)      # save caller's s0
        ...
        lw    ra, 12(sp)     # restore ra
        lw    s0, 8(sp)      # restore s0
        addi  sp, sp, 16     # pop the frame
        ret
\end{lstlisting}

The pairing is the discipline: what goes \texttt{sw} at the top comes
\texttt{lw} at the bottom, in the same offsets, every path out. Miss
either half and the function returns to the wrong place --- the
classic \emph{stack corruption} bug that no error message names.

\begin{example}[follow one frame]\label{ex:frametrace04}
For \texttt{fib(7)}: the outer call pushes a 16-byte frame at
\texttt{sp=0x0ff0}, saves \texttt{ra} and \texttt{s0}, calls
\texttt{fib(6)} --- which pushes another at \texttt{sp=0x0fe0}. At the
deepest, \texttt{fib(1)} is a leaf-ish base case that returns without
recursing. How many frames are live at once? The call tree of fib(7)
reaches depth 7 (fib(7)→fib(6)→…→fib(1)) before the first unwind ---
seven frames, 112 bytes of stack, live simultaneously.
\end{example}

\begin{figure}[t]
\centering
\begin{tikzpicture}[font=\scriptsize, >=latex]
% fib stack frame anatomy
\node[align=left, font=\footnotesize\bfseries] at (4.5,4.7)
  {one \ttfamily fib\rmfamily\ frame (16 B), sp $\to$ low};
\foreach \y/\t/\c/\d in {
  3.8/{ra (12(sp))}/{accentsoft}/{who called \emph{me}},
  3.0/{s0 (8(sp))}/{goodbg}/{caller's saved register},
  2.2/{fib(n-1) (4(sp))}/{graybg}/{parked result},
  1.4/{(unused pad)}/{white}/{16-alignment}}{
  \draw[thick, fill=\c] (0.8,\y-0.3) rectangle (3.2,\y+0.3);
  \node[font=\tiny] at (2.0,\y) {\t};
  \node[right, font=\tiny, text=slate] at (3.4,\y) {\d};
}
\draw[->, thick] (0.5,4.0) -- (0.5,1.0)
  node[midway, left, font=\tiny, rotate=90]{sp grows down};
\node[align=left, text=slate, text width=10cm] at (5.3,0.3)
  {Every non-leaf call pays this: ra saved because the inner
   \ttfamily call\rmfamily\ overwrites it, s0 saved because the
   caller promised to preserve it, a slot for the result the second
   call would clobber. fib(7) builds 13 of these at its deepest.};
\end{tikzpicture}

\caption[The stack frame]{Anatomy of one \texttt{fib} frame. Every
non-leaf call pays this tax: the saved \texttt{ra} keeps the return
path alive, the saved \texttt{s0} honours the callee-saved contract,
and the scratch slot parks \texttt{fib(n-1)} while the second
recursive call runs.}
\label{fig:frame04}
\end{figure}

\section{Caller-saved, callee-saved}
\label{sec:ownership04}
\index{caller-saved}\index{callee-saved}

The rule that confuses everyone on first contact: \texttt{fib} saves
\texttt{ra} --- a \emph{caller}-saved register --- yet the text above
calls that ``the callee's'' job. The resolution is that ownership
follows the \emph{overwrite}, not the register name. \texttt{fib}'s
\texttt{ra} is a value \texttt{fib} still needs; the thing that would
destroy it is \texttt{fib}'s own inner \texttt{call}. So \texttt{fib},
acting as the \emph{caller} of that call, saves it. The same function
plays both roles: it is a callee to whoever called it (so it saves
\texttt{s0} before using it) and a caller to its recursive call (so it
saves \texttt{ra} before that call).

\begin{figure}[t]
\centering
\begin{tikzpicture}[font=\scriptsize, >=latex]
% caller vs callee saved ownership
\draw[thick, fill=accentsoft] (0.4,2.6) rectangle (4.6,4.4);
\node[font=\footnotesize] at (2.5,4.0) {caller-saved (t0--t6, a0--a7, ra)};
\node[align=left, font=\tiny, text width=3.9cm] at (2.5,3.1)
  {the CALLER must save these if it still needs them --- the callee
   may trash them freely};
\draw[thick, fill=goodbg] (5.4,2.6) rectangle (9.6,4.4);
\node[font=\footnotesize] at (7.5,4.0) {callee-saved (s0--s11, sp)};
\node[align=left, font=\tiny, text width=3.9cm] at (7.5,3.1)
  {the CALLEE must restore these before ret --- the caller is allowed
   to believe they're unchanged};
\node[align=left, text=slate, text width=10cm] at (5.0,1.9)
  {One line each way: whoever \emph{loses} the value saves it.
   \ttfamily sw ra,12(sp)\rmfamily\ is the caller... no wait ---
   fib saves its own ra because \emph{fib's} return address dies
   when fib calls itself. Ownership follows the overwrite, not the
   register.};
\end{tikzpicture}

\caption[Caller- vs callee-saved]{Who saves what. Caller-saved
registers are the caller's problem; callee-saved registers are the
callee's problem. One register, \texttt{ra}, sits in the caller-saved
family and is the first thing every non-leaf function saves.}
\label{fig:ownership04}
\end{figure}

\begin{checkpoint}[leaf or not?]\label{cp:leaf04}
A \emph{leaf} function --- one that calls nothing --- needs no frame
at all if it fits in registers and touches no callee-saved ones. Our
\texttt{max} from ch03 was a leaf: \texttt{bge / mv / ret}, three
instructions, zero memory traffic. Predict: does \texttt{empty: ret}
need \texttt{addi sp,sp,-16}? No --- nothing to save.
\end{checkpoint}

\begin{example}[non-leaf: recursive fib]\label{ex:fib04}
\texttt{code/ch04\_fib.S} runs \texttt{fib(7)} on the book's
simulator with a real \texttt{sp} (\texttt{li sp,0x1000}) and a real
stack. The recursion is honest: \texttt{fib} calls itself twice, so
\texttt{ra} must be saved --- and is, at \texttt{12(sp)}. Running it:
\begin{lstlisting}
$ python3 ch03_rv32sim.py ch04_fib.S
13
# 565 instructions retired
\end{lstlisting}
\texttt{fib(7)=13} is correct --- and 565 retired instructions for a
seven-deep recursion is the price tag this chapter is about.
\end{example}

\begin{figure}[t]
\centering
\begin{tikzpicture}[font=\scriptsize, >=latex]
% call tree of fib(7) -- first few nodes
\node[font=\footnotesize\bfseries] at (4.6,4.6) {call tree of \ttfamily fib(7)\rmfamily, depth first};
\foreach \x/\y/\l in {
  5.0/4.0/{f(7)},
  3.2/3.2/{f(6)},
  1.8/2.4/{f(5)},
  0.9/1.6/{f(4)},
  0.4/0.8/{f(3)},
  4.6/3.2/{f(5)},
  6.8/3.2/{f(5)},
  6.8/4.0/{f(6)}}{
  \draw[thick, fill=accentsoft] (\x-0.32,\y-0.2) rectangle (\x+0.32,\y+0.2);
  \node[font=\tiny] at (\x,\y) {\l};}
\draw[->] (4.7,3.8) -- (3.5,3.4);
\draw[->] (5.3,3.8) -- (6.5,3.4);
\draw[->] (3.0,3.0) -- (2.1,2.6);
\draw[->] (1.6,2.2) -- (1.1,1.8);
\draw[->] (0.7,1.4) -- (0.5,1.0);
\node[align=left, text=slate, text width=10cm] at (5.0,-0.3)
  {fib(7) makes 41 calls --- 40 non-leaf frames, each paying the
   16-byte stack tax. Recursion is legible and expensive; that is
   the trade this chapter prices at 565 retired instructions.};
\end{tikzpicture}

\caption[fib(7) call tree]{The call tree of \texttt{fib(7)}. Every
non-leaf node pays a frame; fib(7) makes 41 calls total --- the
recursion is legible, and it is expensive.}
\label{fig:leaftree04}
\end{figure}

\section{What a call costs}
\label{sec:cost04}
\index{calling overhead}

\begin{example}[predict, then measure]\label{ex:cost04}
Prediction first: \texttt{call} = \texttt{jal} = one instruction;
\texttt{ret} = \texttt{jalr} = one instruction. So a call should cost
2 retired instructions beyond the inlined body. Measure it:
\texttt{ch04\_callcost.S} calls \texttt{empty} 100 times;
\texttt{ch04\_inline.S} runs the same loop with the body inlined.
\begin{lstlisting}
$ python3 ch03_rv32sim.py ch04_callcost.S
A
# 410 instructions retired
$ python3 ch03_rv32sim.py ch04_inline.S
B
# 207 instructions retired
\end{lstlisting}
$410-207 = 203$ for 100 calls: just over 2 instructions each ---
the prediction is exact once the loop counter's three instructions
account for the shared work. On a real CPU the frame traffic adds
cache misses on top of this instruction tax.
\end{example}

\begin{figure}[t]
\centering
\begin{tikzpicture}[font=\scriptsize, >=latex]
% call/ret pc flow
\draw[thick, fill=accentsoft] (0.4,2.8) rectangle (3.0,4.0);
\node[align=center, font=\tiny] at (1.7,3.6) {caller};
\node[align=center, font=\tiny] at (1.7,3.1) {\ttfamily call fib\\ra $\leftarrow$ pc+4};
\draw[thick, fill=goodbg] (6.6,2.8) rectangle (9.2,4.0);
\node[align=center, font=\tiny] at (7.9,3.6) {fib};
\node[align=center, font=\tiny] at (7.9,3.1) {\ttfamily ret\\pc $\leftarrow$ ra};
\draw[->, very thick, accent] (3.0,3.4) -- (6.6,3.4)
  node[midway, above, font=\tiny]{pc jumps, ra = return link};
\draw[->, very thick, good] (6.6,3.0) -- (3.0,3.0)
  node[midway, below, font=\tiny]{pc jumps back};
\node[align=left, text=slate, text width=10cm] at (5.0,2.0)
  {Two instructions of mechanism: \ttfamily call\rmfamily\ writes
   ra and jumps; \ttfamily ret\rmfamily\ jumps to ra. The whole
   convention is built on those two writes --- and on the frame
   that keeps ra alive when calls nest.};
\end{tikzpicture}

\caption[call/ret mechanism]{The two-instruction mechanism everything
sits on. \texttt{call} leaves the return link in \texttt{ra};
\texttt{ret} consumes it. Frames exist to keep that link alive across
nesting.}
\label{fig:callret04}
\end{figure}

\section{Reading real compiler output}
\label{sec:readasm04}
\index{compiler output}\index{-O0}\index{-O2}

\begin{example}[-O0 vs -O2 on a switch]\label{ex:switch04}
\texttt{grade(int s)} maps $0\to60, 1\to70, 2\to80, 3\to90$ by
\texttt{switch}. Compiled \texttt{clang -target riscv64 -O0}: 32
instructions --- a 32-byte frame, \texttt{ra}/\texttt{s0} spilled, the
argument stored to the stack, reloaded, \texttt{sext.w}'d, compared
case by case. The same file at \texttt{-O2}: \textbf{nine}
instructions, zero branches ---
\begin{lstlisting}
slli a1, a0, 1        # a1 = 2s
slli a2, a0, 3        # a2 = 8s
sltiu a0, a0, 4       # a0 = (s < 4)
add   a1, a1, a2      # a1 = 10s
addiw a1, a1, 60      # a1 = 60 + 10s
neg   a0, a0          # a0 = -(s<4)
and   a0, a0, a1      # mask: result or 0
ret
\end{lstlisting}
The compiler noticed the cases were linear in \texttt{s} and replaced
the jump table --- and the whole control flow --- with arithmetic and
a sign mask. Prediction failed, delightfully: the real lesson is that
\texttt{-O2} rewrites structure, not just speed.
\end{example}

\begin{figure}[t]
\centering
\begin{tikzpicture}[font=\scriptsize, >=latex]
% O0 vs O2 side by side (real clang riscv64 output counts)
\node[align=left, font=\footnotesize\bfseries] at (4.7,4.6)
  {\ttfamily grade()\rmfamily\ switch: -O0 vs -O2 (real clang 23,
   riscv64 target)};
\draw[thick, fill=graybg] (0.4,1.0) rectangle (4.4,4.2);
\node[font=\footnotesize] at (2.4,3.8) {-O0};
\node[align=left, font=\tiny, text width=3.7cm] at (2.4,2.5)
  {32 instructions:\\frame setup, stores to\\
stack, one branch per case,\\
reload return value, epilogue};
\draw[thick, fill=goodbg] (5.4,1.0) rectangle (9.4,4.2);
\node[font=\footnotesize] at (7.4,3.8) {-O2};
\node[align=left, font=\tiny, text width=3.7cm] at (7.4,2.5)
  {9 instructions, no branches:\\
\texttt{slli, slli, sltiu, add,\\
addiw 60, neg, and, ret} ---\\
the switch became arithmetic};
\node[align=left, text=slate, text width=10cm] at (5.0,0.2)
  {The jump table never appeared: clang proved the four cases were
   linear in n and replaced lookup by computation. Reading
   \ttfamily -O2\rmfamily\ output is how you learn what your C
   really means.};
\end{tikzpicture}

\caption[-O0 vs -O2 codegen]{Real clang~23 output for the same
\texttt{grade()} switch: 32 framed, branching instructions at
\texttt{-O0}; 9 branchless ones at \texttt{-O2}. The jump table never
materialised because the table was linear.}
\label{fig:codegen04}
\end{figure}

\begin{checkpoint}[why the frame at all?]\label{cp:frame04}
At \texttt{-O0} \texttt{grade} spills \texttt{ra} and \texttt{s0} even
though it calls nothing --- the compiler emits the safe frame
unconditionally at that level. At \texttt{-O2} it proves the function
is a leaf and emits no frame at all. That one observation --- leaf
vs.\ non-leaf, proven by the compiler --- is the whole
\texttt{-O0}/\texttt{-O2} gap in miniature.
\end{checkpoint}

\begin{example}[when the compiler deletes your code]\label{ex:volatile04}
A benchmark loop of the form \texttt{sum += i} over \texttt{N}
iterations, at \texttt{-O2}, collapses to a constant --- clang computes
$N(N-1)/2$ at compile time and the loop vanishes. To keep
measurements honest you mark the variable \texttt{volatile} (forces
a real load each iteration) or feed the result into something the
compiler cannot see through. The chapter's own sim-side measurement
needed no such trick: our interpreter does not optimise at all.
\end{example}

\section{PIC in one page}
\label{sec:pic04}
\index{PIC}\index{position-independent code}

Shared libraries cannot know their load address at compile time, so
they are compiled \emph{position-independent}: data access goes
through the GOT, calls go through the PLT --- the same relocation
plumbing ch01 showed. RV32I keeps PIC cheap because \texttt{auipc}
materialises pc-relative addresses in one instruction; our bare-metal
programs sidestep it entirely by linking at absolute addresses.

\section{Common mistakes}
\label{sec:mistakes04}
\begin{itemize}
\item \textbf{Saving \texttt{ra} only sometimes.} A non-leaf function
      must save \texttt{ra} on \emph{every} path --- including error
      returns --- or some returns go to the wrong pc.
\item \textbf{Misaligned \texttt{sp}.} The psABI wants \texttt{sp}
      16-aligned at call sites; \texttt{addi sp,sp,-12} breaks it.
\item \textbf{Using \texttt{s0} without saving.} Callee-saved means
      \emph{the callee} promises; touch \texttt{s0} and you inherit
      the obligation.
\item \textbf{Assuming arguments survive.} \texttt{a0}--\texttt{a7}
      are caller-saved: anything still needed across a call must be
      in an \texttt{s}-register or a frame slot first.
\item \textbf{Trusting \texttt{-O0} codegen to mean anything.}
      \texttt{-O0} writes a frame for a leaf function --- it encodes
      the safe pattern, not the truth about your algorithm.
\item \textbf{Benchmarking at \texttt{-O2} without volatile.}
      The compiler deletes dead loops; what you measured was the
      delete.
\end{itemize}

\section{Summary}
\label{sec:summary04}
\begin{tabularx}{\linewidth}{@{}l X@{}}
\toprule
\textbf{Concept} & \textbf{What it was} \\
\midrule
mechanism & \texttt{call}/\texttt{ret}: two instructions, one register \\
contract & caller-saved vs callee-saved, args \texttt{a0}--\texttt{a7},
 ret \texttt{a0}, \texttt{sp} 16-aligned \\
frame & \texttt{sw ra,12(sp)} pattern; non-leaf only \\
measured & \texttt{fib(7)} = 13 in 565 instrs; call tax = 2 instr \\
compiler & \texttt{-O0} = safe frames; \texttt{-O2} = structure rewritten \\
\bottomrule
\end{tabularx}

\section*{Further reading}
\label{sec:further04}
\begin{itemize}
\item RISC-V psABI specification (the calling convention itself):\\
      \url{https://github.com/riscv-non-isa/riscv-elf-psabi-doc}
\item CS:APP ch3 (procedures, stack discipline):\\
      \url{https://csapp.cs.cmu.edu/}
\item Compiler Explorer (watch -O0/-O2 on your own C):\\
      \url{https://godbolt.org/}
\item The RISC-V Reader, Patterson \& Waterman (calling convention ch):\\
      \url{http://www.riscvbook.com/}
\item GCC -O documentation:\\
      \url{https://gcc.gnu.org/onlinedocs/gcc/Optimize-Options.html}
\end{itemize}

\begin{rulebox}[Watch instead]
\begin{itemize}
\item \textbf{RISC-V General-Purpose Registers: Usage and Calling
      Convention} --- Austin's BSP Lab, 6:35 --- caller-vs-callee-saved
      register roles and why they split that way. Matches
      sec:contract04--sec:ownership04. \url{https://youtu.be/WawdgAEfwvA}
\item \textbf{What Has My Compiler Done for Me Lately?} ---
      Matt Godbolt (CppCon 2017), 75:45 --- Compiler Explorer's author
      reading -O0/-O2 output live. Matches sec:readasm04.
      \url{https://youtu.be/bSkpMdDe4g4}
\end{itemize}
\end{rulebox}

\section*{Notes}
\label{sec:notes04}
\begin{itemize}
\item psABI register roles and \texttt{sp} alignment: riscv-elf-psabi-doc.
\item Frame shape, leaf/non-leaf, -O0 spill behaviour: CS:APP ch3;
      verified against clang 23 riscv64 output captured in
      \texttt{code/ch04\_switch\_O0.s} and \texttt{ch04\_switch\_O2.s}.
\item fib(7)=13 and the 565/410/207 counts: real runs of
      \texttt{ch03\_rv32sim.py}; \texttt{call}/\texttt{j} encoding bug
      fixed in-tree this chapter (0x00800000$\to$0x0080006f).
\end{itemize}

\section{Exercises}
\label{sec:exercises04}

\begin{exercise}[easy]\label{ex:e04a}
Is \texttt{empty: ret} a leaf? Does it need a frame? Justify in one
sentence.
\end{exercise}
\begin{solution}[hint]
Leaf --- no \texttt{call} inside --- and no: nothing it uses can be
clobbered.
\end{solution}

\begin{exercise}[easy]\label{ex:e04b}
Classify \texttt{a0}, \texttt{t2}, \texttt{s3}, \texttt{ra},
\texttt{sp} as caller- or callee-saved.
\end{exercise}
\begin{solution}[hint]
a0,t2,ra caller; s3,sp callee.
\end{solution}

\begin{exercise}[medium]\label{ex:m04a}
\texttt{fib} saves \texttt{ra} at \texttt{12(sp)} but a caller of
\texttt{fib} need not save \texttt{ra}. Who is responsible for which
\texttt{ra}?
\end{exercise}
\begin{solution}[hint]
Every non-leaf saves its OWN ra; the caller never touches the
callee's.
\end{solution}

\begin{exercise}[medium]\label{ex:m04b}
Why is \texttt{fib(n-1)} stored at \texttt{4(sp)} rather than in
\texttt{t0}? What would \texttt{t0} become after the second
\texttt{call}?
\end{exercise}
\begin{solution}[hint]
t0 is caller-saved --- the inner call may trash it; the frame slot
is the safe place.
\end{solution}

\begin{exercise}[medium]\label{ex:m04c}
Predict then check: at \texttt{-O2} would clang still emit a frame
for \texttt{grade}? What does it emit instead of the switch?
\end{exercise}
\begin{solution}[hint]
No frame --- it proves leaf-ness. The switch became arithmetic
(60+10s masked).
\end{solution}

\begin{exercise}[medium]\label{ex:m04d}
\texttt{sum += i} loops disappear at \texttt{-O2}. What are two ways
to keep a benchmark loop alive?
\end{exercise}
\begin{solution}[hint]
volatile on the accumulator, or an opaque sink for the result.
\end{solution}

\begin{exercise}[hard]\label{ex:h04a}
Write a \texttt{sum4(a,b,c,d)} that chains four calls: \texttt{add2}
called twice, results combined. List every register you saved and
justify each.
\end{exercise}
\begin{solution}[hint]
sum4 is non-leaf: save ra + s-registers for the two partial results;
a-regs die on every call.
\end{solution}

\begin{exercise}[hard]\label{ex:h04b}
Design a harness that detects a callee failing to restore
\texttt{s7}: set \texttt{s7}, call the suspect function, check on
return. Write it for the book's simulator.
\end{exercise}
\begin{solution}[hint]
li s7,<sentinel> / call f / compare s7, report mismatch via MMIO.
\end{solution}

\begin{exercise}[hard]\label{ex:h04c}
Predict \texttt{fib(8)}'s instruction count on the sim --- it is
not 565. Then run it and explain the ratio.
\end{exercise}
\begin{solution}[hint]
Calls grow ~1.6x per step (fib calls fib twice); expect $\sim$900,
then measure.
\end{solution}

% chapter end

\begin{exercise}[easy]\label{ex:e04g}
On our ABI, which registers must \texttt{fib} save before touching
them: \texttt{a0} or \texttt{s1}?
\end{exercise}
\begin{solution}[hint]
s1 --- callee-saved; a0 is caller-saved and gets clobbered by the
recursive call anyway.
\end{solution}

\begin{exercise}[easy]\label{ex:e04h}
Why does \texttt{main} need a frame at all if it makes calls but
has no locals?
\end{exercise}
\begin{solution}[hint]
It must save ra --- the call to fib overwrites the return
address main itself needs.
\end{solution}

\begin{exercise}[medium]\label{ex:m04e}
Our leaf-call measurement showed exactly 2 instructions per call.
Reproduce it on a 3-argument function and explain why the count
stays 2 until arg 8.
\end{exercise}
\begin{solution}[hint]
call+ret is the constant cost; a0--a7 are registers --- the 9th
argument spills to the stack.
\end{solution}

\begin{exercise}[medium]\label{ex:m04f}
clang at -O2 turned our switch into arithmetic with no jump
table. Name the transformation and one input shape that makes the
table win anyway.
\end{exercise}
\begin{solution}[hint]
Linearisation/if-conversion; many cases (dense, $>$8) make the
indirect jump cheaper than a compare chain.
\end{solution}

\begin{exercise}[medium]\label{ex:m04g}
The stack grows \emph{down} on RV32I. Show why that makes
\texttt{sp-16} the frame allocation, and what goes wrong if a
function subtracts less than it stores.
\end{exercise}
\begin{solution}[hint]
Lower addresses are free space; under-allocating overwrites the
caller's frame below sp --- silent corruption.
\end{solution}

\begin{exercise}[hard]\label{ex:h04d}
Design the frame for a function with 3 callee-saved regs, one
8-byte local, and calls. Give the exact sp adjustment, each save
slot, and justify the 16-byte alignment.
\end{exercise}
\begin{solution}[hint]
ra + 3 saves + 8B local = 40B; round to 48 for 16B alignment ---
the ABI's sp must stay 16-aligned at call sites.
\end{solution}
